\documentclass[../Main.tex]{subfiles}

\begin{document}

\section{Organization}\label{sec:organization}

Before examining the code itself, it is important to understand how the project is organized. This chapter introduces the file structure of Project Polearm, the purpose of its directories and describes the conventions established throughout development. By grouping related systems, separating shared resources from system-specific implementations and following consistent naming and structural conventions, we sought to make the project understandable and maintainable even as it grew in complexity, as cleanliness is not only important in code, but also where the code resides.

\subsection{Files}\label{sec:files}

Within Project Polearm, the most important and a large portion of files are either of the extensions \code{.gd}, \code{.cs}, \code{.tscn}. These contain the code and scene data needed by Godot to run the game. The code for the game is found in \code{.gd} and \code{.cs} files, whereas \code{.tscn} and occasionally \code{.scn} describe scene files. \code{.tscn} stores its data in a human-readable and editable text format based on the INI format, and \code{.scn} stores its data in a binary format, which is exclusively used for scenes containing large amounts of data. One can also encounter \code{.tres} (text resource) and \code{.res} (resource) files within the project, which denote other data which Godot needs to run the game. \code{.gdshader} files denote shader files in Godot's custom shader language.

One can also find an assortment of files which contain the assets of the game. 3D models are exclusively stored in the GLTF format with the extension \code{.glb}. Textures are most commonly stored as \code{.png} files and \code{.jpeg} files. Godot resource files may also be found here.

Other file types can also be found throughout the project. Unless explicitly denoted otherwise, these contain metadata used by Godot or configuration information for either Godot, .NET and other tools which were used.

\subsection{Conventions}\label{sec:conventions}

The project uses a set of enforced conventions for its code and files to promote navigability, consistency and good practice. They are documented in \code{Common/Templates/Convention/Convention.gd}. This convention fully defines a rigid structure and way of organizing code for all long code files, as well as standardized naming conventions for almost every language feature and pattern used.

Furthermore, we have also orally agreed upon a convention for all file and directory names at the start of development, which simply amounts to naming all files and directories which are not generated automatically in \code{PascalCase}.

Although the naming conventions are enforced throughout most parts of the project, some conventions were enforced with some leniency so as to allow for developer flexibility, as conventions are only good as long as they promote their stated goals, which are navigability, consistency and good practice. The conventions were not enforced for their own sake.

\subsection{Directories}\label{sec:directories}

All directories in the project were named to be as self-explanatory as possible. The entire project encompasses the directories below:

\begin{itemize}
    \item \code{Assets}
    \item \code{addons}
    \item \code{Common}
    \item \code{Components}
    \item \code{Core}
    \item \code{Entities}
    \item \code{Interfaces}
    \item \code{Objects}
    \item \code{Stages}
    \item \code{Systems}
\end{itemize}

We organized our directories as such to best describe their contents as well as their association and dependencies. A directory usually contains source and scene files which belong to the same system, or a grouping of systems which stands on its own pillar within the project. \code{project.godot}, the core project configuration file, sits at the root directory of the project.

Below is a summary of all other directories within Project Polearm and an explanation of their contents in alphabetical order:

\subsubsection{\code{Assets}}\label{sec:assets}

The \code{Assets} directory contains the majority of the game's visual and audio assets, including 3D models, textures, materials, animations, sound effects, music as well as their metadata. The subdirectories under \code{Assets} directly categorize the type of assets contained within. Within those, assets are either sorted by their functionality, their associated system and/or their origin, allowing them to be located and referenced relatively easily.

It is important to note that the assets in this directory consist of external assets. That means that they do not directly influence anything about the execution of the game. Other assets, specifically game data, like animation frame data or asset datasets, which are crucial to the behavior of the game, are stored alongside their respective systems and entities.

The directories are named to be self-explanatory and will therefore not be explained further.

\subsubsection{\code{addons}}\label{sec:addons}

This directory contains all Godot plugins used throughout the game's development. As mentioned, only \code{PhantomCamera}, found in the folder of the same name, is necessary for the game's execution, but other extensions were kept for completeness' sake.

The plugins were installed through Godot's built-in Asset Library, which downloads them into this specific directory. The name of the directory uses Godot's own convention, which we did not enforce our own onto, as that may cause issues and inconveniences. Therefore it is the only directory in the project, created by us, which does not conform to our naming convention.

These are all addons used in the project and how we utilized them in Project Polearm:

\begin{table}[H]
    \centering
    \small
    \renewcommand{\arraystretch}{1.3}
    \begin{tabular}{|>{\raggedright\arraybackslash}p{3.0cm}|>{\raggedright\arraybackslash}p{3.0cm}|>{\raggedright\arraybackslash}p{5.0cm}|>{\raggedright\arraybackslash}p{3.2cm}|}
        \hline
        \textbf{Extension / directory name} & \textbf{Functionality} & \textbf{Reason for utilization} & \textbf{Runtime dependency} \\
        \hline
        \textbf{\code{anim\_player\_-}} \linebreak \textbf{\code{refactor}} &
        \begin{minipage}[t]{\linewidth}\raggedright
            Edit Godot scene paths within animations\vspace{0.4em}
        \end{minipage}
        &
        \begin{minipage}[t]{\linewidth}\raggedright
            Godot usually imports models grouped alongside their skeleton and animations, which we separated. Therefore, all animation paths become invalid and must be reset\vspace{0.4em}
        \end{minipage}
        &
        \begin{minipage}[t]{\linewidth}\raggedright
            No, editor plugin\vspace{0.4em}
        \end{minipage}
        \\
        \hline
        \textbf{\code{debug\_draw\_3d}} &
        \begin{minipage}[t]{\linewidth}\raggedright
            Draw simple, geometrical shapes for debugging\vspace{0.4em}
        \end{minipage}
        &
        \begin{minipage}[t]{\linewidth}\raggedright
            This plugin was used during development to display 3D debugging information, mainly for the inspection of vectors\vspace{0.4em}
        \end{minipage}
        &
        \begin{minipage}[t]{\linewidth}\raggedright
            Runtime plugin, no release build dependency\vspace{0.4em}
        \end{minipage}
        \\
        \hline
        \textbf{\code{phantom\_camera}} &
        \begin{minipage}[t]{\linewidth}\raggedright
            Ready-to-go camera library for Godot\vspace{0.4em}
        \end{minipage}
        &
        \begin{minipage}[t]{\linewidth}\raggedright
            Godot cameras only implement basic features, with the rest having to be implemented by the developer. Phantom Camera implements functionality useful for most games, like following, smooth camera movement, camera collision etc.\vspace{0.4em}
        \end{minipage}
        &
        \begin{minipage}[t]{\linewidth}\raggedright
            Runtime plugin, release build dependency\vspace{0.4em}
        \end{minipage}
        \\
        \hline
    \end{tabular}
    \caption{Godot plugins used in Project Polearm.}
    \label{tab:addons}
\end{table}

\subsubsection{\code{Common}}\label{sec:common}

\code{Common} is the place for everything which does not belong to a single system, but is needed by several of them. What separates it from \code{Components} is that nothing in here is meant to be dropped into a scene and configured on its own. The code is either supposed to be extended, called from elsewhere, or simply a utility.

{\small
\renewcommand{\arraystretch}{1.3}
\begin{longtable}{|>{\raggedright\arraybackslash}p{2.2cm}|>{\raggedright\arraybackslash}p{8.0cm}|>{\raggedright\arraybackslash}p{4.0cm}|}
    \caption{Contents of the Common directory.}\label{tab:common} \\
    \hline
    \textbf{Directory} & \textbf{Contents} & \textbf{See also} \\
    \hline
    \endfirsthead

    \hline
    \textbf{Directory} & \textbf{Contents} & \textbf{See also} \\
    \hline

    \endhead
    \multicolumn{3}{r}{\textit{Continued on next page}} \\
    \endfoot

    \hline
    \endlastfoot

    \code{StateMachine} &
    \code{BasicStateMachine} and \code{BasicState}, the generic finite state machine and its base state. All state machines in the game, which include entities, menus and the application itself, either inherit directly from the pair or use a similar implementation if that is not possible &
    Section~\ref{sec:state-machine} \\
    \hline

    \code{Utility} &
    \code{PPUtil}, a collection of constants, enumerations and math helpers, e.g. for rotating a vector onto a plane. As the project uses both GDScript and C\#, it exists in two versions, \code{Utility.gd} and \code{Utility.cs}, for ease of use in both languages &
    - \\
    \hline

    \code{SceneCamera} &
    \code{SceneCamera}, the main camera of the game. It contains the host of \textbf{PhantomCamera}, whose transformation is decided by sub-cameras such as the player camera, as well as a post-processing layer. It is kept as its own scene to allow for swappable cameras within the \code{Visuals} system &
    Sections~\ref{sec:entities}, \ref{sec:systems} \\
    \hline

    \code{Templates} &
    \code{Convention}, where the coding conventions of the project are written down, and \code{PlayerAction}, a script template for new player states &
    Section~\ref{sec:conventions} \\
    \hline

    \code{Tools} &
    Scripts and configuration files which only serve the developer inside the editor &
    Chapter~\ref{sec:tooling-and-developer-experience} \\
    \hline
\end{longtable}
}

\subsubsection{\code{Components}}\label{sec:components}

Components are a cornerstone of Project Polearm's architecture and a result of our development philosophies (see section~\ref{sec:configure-everything-code-nothing-our}). Each one is a small, self-contained building block, which can be added to a scene, configured through its exported variables and made to work with minimal or no extra coding. Components are almost entirely self-contained and their core implementations know about no other big systems.

{\small
\renewcommand{\arraystretch}{1.3}
\begin{longtable}{|>{\raggedright\arraybackslash}p{0.13\linewidth}|>{\raggedright\arraybackslash}p{0.51\linewidth}|>{\raggedright\arraybackslash}p{0.26\linewidth}|}
    \caption{Components of Project Polearm.}\label{tab:components} \\
    \hline
    \textbf{Directory} & \textbf{Functionality} & \textbf{Used by} \\
    \hline
    \endfirsthead

    \hline
    \textbf{Directory} & \textbf{Functionality} & \textbf{Used by} \\
    \hline
    \endhead

    \hline
    \multicolumn{3}{r}{\textit{Continued on next page}} \\
    \endfoot

    \hline
    \endlastfoot

    \code{Collision} &
    \code{Hitbox} and \code{Hurtbox} for dealing and receiving damage respectively. \code{HitboxData} is the collision payload and the bridge between the two components (see section~\ref{sec:from-attacking-to-damage}) &
    Weapons, entities, kill zones \\
    \hline

    \code{EventTrigger} &
    \code{EventTrigger}, an area which triggers when it is entered &
    \code{EncounterTrigger} \\
    \hline

    \code{FloatingBody} &
    A custom character physics body written in C\#. It is split into three files: \code{FloatingBody.cs} for the configuration, \code{FloatingBodyAPI.cs} for public API methods, and \code{FloatingBodyPhysics.cs} for the physics algorithm implementation (see sections~\ref{sec:floatingbody} and~\ref{sec:a-dream-of-stepping-up}) &
    All entities \\
    \hline

    \code{Interaction} &
    \code{Interactor} and \code{Interactable}, a pair for interactions with objects. Currently unused, as the demo does not implement interactables &
    - \\
    \hline

    \code{LedgeStopper} &
    \code{LedgeStopper} and \code{LedgeStopperBlock}, an activatable component which can stop objects from going over ledges &
    The player \\
    \hline

    \code{Miscellaneous} &
    \code{RemoteChild}, a small node which copies the position of another node without having to be its scene child &
    \code{PlayerCameraManager} \\
    \hline

    \code{Numeric-} \linebreak \code{Resource} &
    \code{ValueResource} and \code{IntegerResource}, which contain a value \(\in[0, max]\) and an interaction API &
    Health, energy and other "resource-like" values \\
    \hline

    \code{ValueBar} &
    \code{ValueBar}, a bar which displays the state of a \code{ValueResource} &
    Health and energy bars \\
\end{longtable}
}

\subsubsection{\code{Core}}\label{sec:core}

\code{Core} is the heart of the game and contains everything that controls the macro-level state of the application. This is implemented as two state machines nested in each other (see sections~\ref{sec:architectural-overview} and~\ref{sec:boot-sequence}).

{\small
\renewcommand{\arraystretch}{1.3}
\begin{longtable}{|>{\raggedright\arraybackslash}p{0.18\linewidth}|>{\raggedright\arraybackslash}p{0.72\linewidth}|}
    \caption{Contents of the Core directory.}\label{tab:core} \\
    \hline
    \textbf{Directory} & \textbf{Contents} \\
    \hline
    \endfirsthead

    \hline
    \textbf{Directory} & \textbf{Contents} \\
    \hline
    \endhead

    \hline
    \multicolumn{2}{r}{\textit{Continued on next page}} \\
    \endfoot

    \hline
    \endlastfoot

    \code{AppStateMachine} &
    The shared base of the core state machines. \code{AppStateMachine} extends the \code{BasicStateMachine} of \code{Common} by implementing a loading state for dynamic state loading. \code{Loading} contains this loading state and the base of all loading screens \\
    \hline

    \code{Main} &
    The outer state machine. \code{Main.tscn} is, as the name suggests, the game's main scene. \code{Boot} contains the splash screen and applies the settings, \code{Loading} the loading screen of \code{Main}, and \code{States} the states of the application: the main menu, the \code{GameState} which loads the \code{Game} state machine, and quitting. \code{ThreadLoadedScene.gd} is the base of all states whose scene is loaded on a background thread \\
    \hline

    \code{Game} &
    The inner state machine, which is loaded when the player enters the game. \code{World} loads the stage and spawns the player, \code{States} contains the pause and death menu states, and \code{Loading} contains the loading screen of \code{Game} \\
    \hline
\end{longtable}
}

Note that two different files are named \code{GameState.gd}. \code{Core/Main/States/GameState.gd} is the state of \code{Main} which loads and enters the game, whereas \code{Core/Game/GameState.gd} is the base class of all states of \code{Game}.

\subsubsection{\code{Entities}}\label{sec:entities}

Entities are the characters of the game, which are the player and all enemies. Instead of giving each of them its own implementation, they all inherit from a shared base and extend it with what makes them different (see section~\ref{sec:architecture-entities}):

\begin{verbatim}
Entity
--> Player
--> BasicEnemy
    --> Solider (variant: HeavySolider)
    --> Archer (variants: ChainArcher, HeavyArcher)
\end{verbatim}

\code{Entity} holds everything every character needs: the action manager (a state machine), the animation and frame data implementation, damage logic, the combo system and other resources like health. \code{Player} extends the basic entity by input handling, a camera and advanced combat, while \code{BasicEnemy} adds an AI controller. The \code{Solider} and \code{Archer}, which extend from \code{BasicEnemy}, are the two basic enemy types in the demo, which are a melee fighter and a ranged fighter respectively. The variants are further enemies derived from them.

Every entity directory is organized in a similar fashion, although not every entity needs every system and therefore every directory:

{\small
\renewcommand{\arraystretch}{1.3}
\begin{longtable}{|>{\raggedright\arraybackslash}p{0.12\linewidth}|>{\raggedright\arraybackslash}p{0.61\linewidth}|>{\raggedright\arraybackslash}p{0.17\linewidth}|}
    \caption{Shared directory structure of the entities.}\label{tab:entities} \\
    \hline
    \textbf{Directory} & \textbf{Contents} & \textbf{See also} \\
    \hline
    \endfirsthead

    \hline
    \textbf{Directory} & \textbf{Contents} & \textbf{See also} \\
    \hline
    \endhead

    \hline
    \multicolumn{3}{r}{\textit{Continued on next page}} \\
    \endfoot

    \hline
    \endlastfoot

    \code{Actions} &
    The action manager, i.e. the state machine of the entity, and in \code{States}, one script per action the entity can perform &
    Sections~\ref{sec:architecture-entities}, \ref{sec:entity-action-selection} \\
    \hline

    \code{Animation} &
    The animator of the entity and, occasionally, its frame data manager &
    Section~\ref{sec:frame-data} \\
    \hline

    \code{Controller} &
    This directory exists for enemies only, as it contains their AI/controller, which turns the role assigned by the \code{EnemyManager} into concrete intents &
    Section~\ref{sec:enemy-role-assignment} \\
    \hline

    \code{Data} &
    Configuration scripts of the entity, such as \code{ActionData}, which configures an entity action, and \code{Resources}, its resource container &
    - \\
    \hline

    \code{Resources} &
    The configured resource files, such as frame data, animation blend trees and combos &
    Section~\ref{sec:frame-data} \\
    \hline

    \code{Variants} &
    Also only for enemies, this directory contains the enemy variants &
    - \\
    \hline
\end{longtable}
}

The difference between \code{Data} and \code{Resources} is worth remembering, as it applies to the entire project: \code{Data} contains the scripts which define \textbf{what} can be configured, whereas \code{Resources} contains the actual configuration.

Some directories only exist for specific entities:

\begin{itemize}
    \item \code{Entity/Collisions} contains \code{EntityHurtbox}, the hurtbox which forwards hits to the entity (see section~\ref{sec:from-attacking-to-damage}).
    \item \code{Entity/Combat} contains the \code{EntityCombatManager} as well as the structure of a combo in \code{Combo}. The configured combos of the player can be found in \code{Player/Resources/Combos}.
    \item \code{Entity/Utility} contains the \code{EntityAnimationHelper} (see section~\ref{sec:entityanimationhelper}).
    \item \code{Player/Input} contains the \code{PlayerInputManager}, the controller of the player, which converts input events into intents.
    \item \code{Player/Camera} contains the \code{PlayerCameraManager}, a C\# script controlling the player's camera, which builds upon \textbf{PhantomCamera}.
    \item \code{Player/Actions/States} contains the 26 states of the player. The directory \code{BasicActions} holds the base classes shared by several states, such as jumping, dashing and stopping, \code{TransitionAction} holds short states which chain into other states, and \code{Physics} holds \code{PlayerPhysics}, the physics configuration shared across states (see section~\ref{sec:to-configure-everything}) as well as ledge grabbing physics.
\end{itemize}

\subsubsection{\code{Interfaces}}\label{sec:interfaces}

\code{Interfaces} stores the menus of the game. Menus, like the core game and entities, are state machines, but they are built on top of the \code{InterfaceStateMachine}, a state machine implementation of their own with a small input manager to boot. It is found in a directory of the same name. Interfaces use a special state machine as they have to inherit from a special \code{Node} type due to how Godot handles interfaces, although their core logic remains the same as that of the \code{BasicStateMachine}.

The menus themselves are found in the directory \code{Menus}:

{\small
\renewcommand{\arraystretch}{1.3}
\begin{longtable}{|>{\raggedright\arraybackslash}p{0.15\linewidth}|>{\raggedright\arraybackslash}p{0.75\linewidth}|}
    \caption{Menus of the game.}\label{tab:menus} \\
    \hline
    \textbf{Directory} & \textbf{Contents} \\
    \hline
    \endfirsthead

    \hline
    \textbf{Directory} & \textbf{Contents} \\
    \hline
    \endhead

    \hline
    \multicolumn{2}{r}{\textit{Continued on next page}} \\
    \endfoot

    \hline
    \endlastfoot

    \code{MainMenu} &
    The main menu, consisting of its screen (\code{MainMenuScreen}), the quit confirmation (\code{QuitConfirmMenu}) and the credits screen \\
    \hline

    \code{CreditsMenu} &
    The credits menu (\code{CreditsMenu}) \\
    \hline

    \code{PauseMenu} &
    The pause menu (\code{PauseScreen}) \\
    \hline

    \code{DeathMenu} &
    The death menu (\code{DeathScreen}), from which the player can respawn or return to the main menu \\
    \hline

    \code{OptionsMenu} &
    The options, shared by the main and pause menus, with their pages \code{GameOptions}, \code{VideoOptions}, \code{AudioOptions} and \code{MappingOptions}. The latter exists in a keyboard and a controller version and uses an \code{ActionMappingButton} for every action that can be remapped \\
    \hline

    \code{Button} &
    \code{GameMenuButton}, the button shared by all menus \\
    \hline
\end{longtable}
}

All menus use the single theme found in \code{Themes}. Furthermore, the \code{Language} directory contains the \code{Project-} \linebreak \code{PolearmLanguages.csv} table, which contains all localization data for the game.

\subsubsection{\code{Objects}}\label{sec:objects}

Objects are, as their name might suggest, inanimate things placed into or spawned in the world. Objects may be thought of as static components of the game world.

{\small
\renewcommand{\arraystretch}{1.3}
\begin{longtable}{|>{\raggedright\arraybackslash}p{0.16\linewidth}|>{\raggedright\arraybackslash}p{0.6\linewidth}|>{\raggedright\arraybackslash}p{0.14\linewidth}|}
    \caption{Contents of the Objects directory.}\label{tab:objects} \\
    \hline
    \textbf{Directory} & \textbf{Contents} & \textbf{See also} \\
    \hline
    \endfirsthead

    \hline
    \textbf{Directory} & \textbf{Contents} & \textbf{See also} \\
    \hline
    \endhead

    \hline
    \multicolumn{3}{r}{\textit{Continued on next page}} \\
    \endfoot

    \hline
    \endlastfoot

    \code{Weapons} &
    \code{Weapon}, the base of melee weapons carried by entities. The weapon of the player can be found in \code{PlayerWeapon}, whereas the \code{Hammer} and \code{Scimitar} of the \code{Solider} enemies are in \code{SoliderWeapon} &
    Section~\ref{sec:from-attacking-to-damage} \\
    \hline

    \code{Weapons/} \linebreak \code{Projectile} &
    \code{Projectile}, a basic projectile object, and \code{ProjectilePhysics}, its physics payload. Its two children \code{Arrow} and \code{HomingMissile} are also found here. Unlike the melee weapons, projectiles are game objects, which are managed by the \code{GameObjectManager} &
    Section~\ref{sec:gameobjectmanager} \\
    \hline

    \code{Door} &
    \code{Door} is an openable and closable, force-field-like door used for level progression &
    Section~\ref{sec:encounter-and-wave} \\
    \hline

    \code{KillZone} &
    The \code{KillZone}, a \code{Hitbox} configured to damage everything inside it by default &
    Section~\ref{sec:from-attacking-to-damage} \\
    \hline
\end{longtable}
}

\subsubsection{\code{Stages}}\label{sec:stages}

A stage is a scene which the game can load as the environment in which the game plays out, consisting of the level geometry, doors and enemy encounters (see section~\ref{sec:architectural-overview}). The base script of all stages is \code{Stage.gd}, located directly in this directory. It defines the stage types and holds the stage configuration data.

{\small
\renewcommand{\arraystretch}{1.3}
\begin{longtable}{|>{\raggedright\arraybackslash}p{0.16\linewidth}|>{\raggedright\arraybackslash}p{0.57\linewidth}|>{\raggedright\arraybackslash}p{0.16\linewidth}|}
    \caption{Contents of the Stages directory.}\label{tab:stages} \\
    \hline
    \textbf{Directory} & \textbf{Contents} & \textbf{See also} \\
    \hline
    \endfirsthead

    \hline
    \textbf{Directory} & \textbf{Contents} & \textbf{See also} \\
    \hline
    \endhead

    \hline
    \multicolumn{3}{r}{\textit{Continued on next page}} \\
    \endfoot

    \hline
    \endlastfoot

    \code{Data} &
    \code{StageData}, which holds the player spawn transform, whereas \code{StageDataset} contains the paths of all stages mapped to their identifiers &
    Section~\ref{sec:dataset-pattern} \\
    \hline

    \code{Levels} &
    \code{Level1}, the demo level, and \code{Encounter.gd}, an encounter script which opens and closes configured doors &
    Sections~\ref{sec:encounter-and-wave}, \ref{sec:encountertrigger} \\
    \hline

    \code{Documentation} &
    Stages which were made for development purposes rather than for the player. The \code{Gym} is an open testing area used during development, with its earlier version and obstacles in \code{Gym/Old}. The \code{Museum} is a test stage built with TrenchBroom &
    Section~\ref{sec:debugging-and-testing} \\
    \hline
\end{longtable}
}

\subsubsection{\code{Systems}}\label{sec:systems}

These are the core systems of Project Polearm, which exist once per game and provide services which other parts of the game depend on. They are singletons managed by \code{Globals.gd}, which lies directly in the root of this directory. See section~\ref{sec:global-state} for more details about their implementation.

{\small
\renewcommand{\arraystretch}{1.3}
\begin{longtable}{|>{\raggedright\arraybackslash}p{0.13\linewidth}|>{\raggedright\arraybackslash}p{0.51\linewidth}|>{\raggedright\arraybackslash}p{0.26\linewidth}|}
    \caption{Global systems of the game.}\label{tab:systems} \\
    \hline
    \textbf{Directory} & \textbf{Contents} & \textbf{See also} \\
    \hline
    \endfirsthead

    \hline
    \textbf{Directory} & \textbf{Contents} & \textbf{See also} \\
    \hline
    \endhead

    \hline
    \multicolumn{3}{r}{\textit{Continued on next page}} \\
    \endfoot

    \hline
    \endlastfoot

    \code{Audio} &
    The \code{AudioManager} for music and, in \code{SFX}, the \code{SFXManager} for 3D sound effects. \code{Data} defines the sound data and datasets, and \code{Resources} contains the configured datasets &
    Section~\ref{sec:architectural-overview} \\
    \hline

    \code{EnemyManager} &
    The \code{EnemyManager}, which ticks, spawns and assigns roles to all enemies. The \code{Encounter} folder contains the \code{Encounter} and \code{Wave} structures, the \code{EncounterTrigger} and the base of all encounter scripts. \code{Perception} contains \code{EnemyPerception}, a subsystem of \code{EnemyManager} which manages simple enemy vision &
    Sections~\ref{sec:enemy-role-assignment}, \ref{sec:encounter-and-wave}, \ref{sec:encountertrigger}, \ref{sec:the-troubles-of-managing-enemies} \\
    \hline

    \code{GameObject} &
    \code{GameObject}, the \code{GameObjectManager} and the \code{GameObjectPool}. \code{Data} defines the dataset of all game objects, and \code{Resources} contains its configured instance &
    Section~\ref{sec:gameobjectmanager} \\
    \hline

    \code{Settings} &
    The \code{SettingsManager}, which manages, stores and applies the game's settings &
    Section~\ref{sec:architectural-overview} \\
    \hline

    \code{Visuals} &
    \code{Visuals}, which holds the \code{SceneCamera} and the world environment of the game &
    Section~\ref{sec:architectural-overview} \\
    \hline
\end{longtable}
}

Besides these systems, \code{Globals} also holds a reference to the \code{World} of \code{Core}, so that other parts of the game can reach the game world.

\newpage

\end{document}
